
本章建立 \BookRecordIndex{T0003}$J$-全纯曲线的基本性质。其中包括\BookRecordIndex{T0249}能量恒等式，它为第 4 章的紧致性论证奠定
基础；还包括命题 2.5.1 对\BookRecordIndex{T0091}单纯曲线的刻画，这是第 3 章\BookRecordIndex{T0166}横截性方法的一个关键要素。另一个
要点是\BookRecordIndex{T0333}唯一延拓定理：在一点具有相同无穷阶 jet 的两条曲线必然重合。我们将利用 \BookRecordIndex{T0013}Carleman
相似原理给出自足的证明。

本章包含许多基础细节。为求完整而收入这些内容，但理解后续章节并不都需要它们。特别地，
尽管最终目标是理解 $C^\infty$ 曲线，第 3、6 章的论证要使用\BookRecordIndex{T0053}巴拿赫流形，而这些流形的
元素是光滑性较低的曲线。因此，本章许多结果是对 $C^\ell$ \BookRecordIndex{T0088}殆复结构证明的。（依照 Floer
的一个想法，也可以在光滑范畴中发展整套理论，注 3.2.7 对此有简要说明。）初学者宜阅读
各节开头主要定理的陈述及较容易的证明，把技术性更强的部分留到用得上时再读。

\section{\BookRecordIndex{T0088}殆复结构}\label{sec:2-1}
\zhindex{daifujiegou}{殆复结构}

设 $(M,\omega)$ 是\BookRecordIndex{T0358}辛流形。\BookRecordIndex{T0088}殆复结构 $J:TM\to TM$ 若满足
\[
  v\ne0\quad\Longrightarrow\quad\omega(v,Jv)>0
\]
对每个 $v\in TM$ 都成立，就称为被 $\omega$ \BookRecordIndex{T0372}驯服的；若它被 $\omega$ \BookRecordIndex{T0372}驯服，并且对
每个 $x\in M$ 以及 $v,w\in T_xM$ 都有
\[
  \omega(Jv,Jw)=\omega(v,w),
\]
则称它与 $\omega$ \BookRecordIndex{T0348}相容。每个被 $\omega$ \BookRecordIndex{T0372}驯服的\BookRecordIndex{T0088}殆复结构都确定一个黎曼度量
\begin{equation}\label{eq:2-1-1}
  g_J(v,w):=\langle v,w\rangle_J
  :=\frac12\bigl(\omega(v,Jw)+\omega(w,Jv)\bigr).
\end{equation}
在 $\omega$-\BookRecordIndex{T0348}相容情形，这个度量可写成
\[
  \langle\,\cdot,\cdot\,\rangle_J=\omega(\,\cdot,J\cdot\,).
\]
无论哪种情形，自同构 $J$ 都是 $TM$ 上的斜自伴等距映射。在黎曼曲面
$(\Sigma,j_\Sigma,d\operatorname{vol}_\Sigma)$ 上，体积形式（或\BookRecordIndex{T0363}辛形式）
$d\operatorname{vol}_\Sigma$ 与复结构 $j_\Sigma$ 一起确定 $\Sigma$ 上的度量。反过来，
若 $\Sigma$ 已定向，则 $\Sigma$ 上每个度量都同时确定一个体积形式和一个复结构。

本书始终以 $\mathcal J(M,\omega)$ 表示 $M$ 上与 $\omega$ \BookRecordIndex{T0348}相容的\BookRecordIndex{T0088}殆复结构组成的空间，
以 $\mathcal J_\tau(M,\omega)$ 表示被 $\omega$ \BookRecordIndex{T0372}驯服的\BookRecordIndex{T0088}殆复结构组成的空间。因此
$\mathcal J(M,\omega)\subset\mathcal J_\tau(M,\omega)$；空间
$\mathcal J_\tau(M,\omega)$ 在 $M$ 上所有\BookRecordIndex{T0088}殆复结构组成的空间 $\mathcal J(M)$ 中是开的，
并且前两个空间都是可缩的（比较 \srcite{277}）。上述定义对任意非退化 $2$-形式
$\omega$ 都适用，无论它是否闭合。不过，除非明确说明相反情形，我们总假设 $\omega$ 是
闭的。

下文以 $N=N_J\in\Omega^2(M,TM)$ 表示 $J$ 的 \BookRecordIndex{T0042}Nijenhuis 张量。它由
\footnote{为保持一致，我们沿用 McDuff--Salamon \srcite{277} 的符号约定，把两个向量场
$X,Y\in\Vect(M)$ 的李括号定义为
\[
  [X,Y]:=\nabla_YX-\nabla_XY=-\mathcal L_XY.
\]
这里 $\nabla$ 是 $M$ 上任一黎曼度量的 \BookRecordIndex{T0034}Levi--Civita 联络。理由见
\srcite{277}，注 3.3。}
\[
  N(X,Y)=[X,Y]+J[JX,Y]+J[X,JY]-[JX,JY]
\]
给出，其中 $X,Y\in\Vect(M)$；它衡量 $J$ 的不可积性。更准确地，$N_J=0$ 当且仅当
$J$ 可积，即 $M$ 有一个坐标卡图册，其转移函数都是 $J$-全纯的。若此外 $J$ 与
$\omega$ \BookRecordIndex{T0348}相容，三元组 $(M,J,g_J)$ 就是\BookRecordIndex{T0207}凯勒流形。一般情形的进一步讨论见
\srcite{277} 第 4 章及其中所列参考文献。
\footnote{还要指出，Hill--Taylor \srcite{173} 的一篇较新论文把
Newlander--Nirenberg 可积性定理推广到了\BookRecordIndex{T0408}正则性略低于 Lipschitz 的\BookRecordIndex{T0088}殆复结构。}
特别要注意，当 $M$ 为 $2$ 维时 $N_J$ 总是消失，所以黎曼曲面上的任意\BookRecordIndex{T0088}殆复结构都可积。
我们为这一特殊情形给出两个证明：一个是第 C.5 节中的推论 C.5.3，另一个在第 E.3 节。

设 $(M,J)$ 是殆复流形，$\omega$ 是 $M$ 上一个不一定闭的 $2$-形式，并且
$\omega(\cdot,J\cdot)$ 是黎曼度量。令 $\nabla$ 表示该度量的 \BookRecordIndex{T0034}Levi--Civita 联络。
则条件 $\nabla J=0$ 刻画凯勒情形，即
\[
  \nabla J=0\quad\Longleftrightarrow\quad d\omega=0\ \text{且}\ N_J=0.
\]
特别地，如果 $(M,\omega)$ 是\BookRecordIndex{T0358}辛流形，而且 $J\in\mathcal J(M,\omega)$ 与 $\omega$
\BookRecordIndex{T0348}相容，那么 $J$ 在 \BookRecordIndex{T0034}Levi--Civita 联络 $\nabla$ 下通常并不平行。另一方面，修正后的联络
\[
  \widetilde\nabla_XY:=\nabla_XY-\frac12J(\nabla_XJ)Y
\]
总保持 $J$ 和度量（即它是\BookRecordIndex{T0052}埃尔米特联络），不过只有在可积情形下它才无挠。事实上，直接
计算可把它的\BookRecordIndex{T0248}挠率 $\widetilde T(X,Y)$ 表示成 \BookRecordIndex{T0042}Nijenhuis 张量：
\[
  \widetilde T(X,Y)
  :=\widetilde\nabla_YX-\widetilde\nabla_XY-[X,Y]
  =-\frac14N(X,Y).
\]
在 $\omega$-\BookRecordIndex{T0372}驯服情形下，$\widetilde\nabla$ 仍是\BookRecordIndex{T0052}埃尔米特联络，但它的\BookRecordIndex{T0248}挠率不再是
\BookRecordIndex{T0042}Nijenhuis 张量。此时还存在另一个复联络 $\widehat\nabla$，其\BookRecordIndex{T0248}挠率是 \BookRecordIndex{T0042}Nijenhuis 张量，
但未必保持度量（见附录 C）。注意，一般而言，$TM$ 上未必存在既无挠又保持 $J$ 的联络。

\begin{exercise}\label{ex:2-1-1}
设 $J$ 是 $M$ 上的\BookRecordIndex{T0088}殆复结构，$\nabla$ 是 $M$ 上任一黎曼度量的 \BookRecordIndex{T0034}Levi--Civita 联络。
证明 $J$ 沿向量场 $X$ 的李导数为
\[
  (\mathcal L_XJ)Y=(\nabla_XJ)Y-\nabla_{JY}X+J\nabla_YX.
\]
\end{exercise}
